\documentclass[11pt]{article}

\usepackage[margin=1in]{geometry}
\usepackage{enumitem}
\usepackage{listings}
\usepackage{hyperref}

\lstdefinestyle{shell}{
  basicstyle=\ttfamily\small,
  columns=fullflexible,
  keepspaces=true,
  frame=single,
  breaklines=true,
  showstringspaces=false
}

\title{Yalnix Checkpoint 3 Test Notes}
\author{COSC 58}
\date{\today}

\begin{document}
\maketitle

\section{Test Commands}

Run from the Thayer source directory:

\begin{lstlisting}[style=shell]
make clean
make
timeout 5 ./yalnix -W
timeout 3 ./yalnix init
\end{lstlisting}

\section{Expected Observations}

\begin{itemize}[noitemsep]
  \item \texttt{make} builds both the kernel executable \texttt{yalnix} and the userland executable \texttt{init}.
  \item \texttt{KernelStart} creates idle and init processes, switches to init's Region 1, and \texttt{LoadProgram} loads \texttt{init}.
  \item Userland \texttt{init} prints its PID from \texttt{GetPid}.
  \item The user library issues an early \texttt{Brk}; the kernel maps heap pages and returns success.
  \item Clock traps alternate execution between \texttt{init} and idle while \texttt{init} pauses.
  \item Every fifth init loop calls \texttt{Delay(2)}; init blocks, idle runs for clock ticks, then init resumes.
\end{itemize}

\section{Run Evidence}

The explicit \texttt{./yalnix init} run checks that \texttt{cmd\_args[0]} selects the init executable. The no-argument run checks the default init name.

Representative trace lines from the Checkpoint 3 smoke test:

\begin{lstlisting}[style=shell]
Yalnix    KernelStart: leaving for init pid 1
Hardware  | Syscall trap Brk
Hardware  | Syscall trap GetPid
User Prog init: pid 1 starting
User Prog init: loop 0 pid 1
Yalnix    KernelStart: leaving for idle pid 0
Yalnix    DoIdle
User Prog init: loop 1 pid 1
Yalnix    DoIdle
User Prog init: loop 4 pid 1
Hardware  | Syscall trap Delay
Yalnix    DoIdle
Yalnix    DoIdle
User Prog init: loop 5 pid 1
\end{lstlisting}

\end{document}
